\documentclass[../Main.tex]{subfiles}
\begin{document}
  \chapter{2020-2021学年微积分（一）（下）期末考试}

  \section{单项选择题(每小题 3 分，共 18 分)}
  \begin{enumerate}
    \item 微分方程 \( y'' - 3y' + 2y = 2x - 2e^x \) 的特解 \( y^* \) 的形式是（\qquad）.

      \begin{tasks}[label=\textbf{\Alph*.}, label-width=1.5em, column-sep=1em](2)
        \task \( y^* = (Ax + B)e^x \)
        \task \( y^* = x(Ax + B)e^x \)
        \task \( y^* = Ax + B + Ce^x \)
        \task \( y^* = Ax + B + Cxe^x \)
      \end{tasks}

    \item 设曲线 \( \begin{cases} x^2 + y^2 + z^2 = 2 \\ x + y + z = 0 \end{cases} \)，点 \( M(1, -1, 0) \)，则在点 \( M \) 处下列说法\(\textbf{不正确}\)的是（\qquad）.

      \begin{tasks}[label=\textbf{\Alph*.}, label-width=1.5em, column-sep=1em](2)
        \task 切矢量为 \(\{-2, -2, 4\}\)
        \task 切矢量为 \(\{-2, 2, 4\}\)
        \task 切线方程为 \( x - 1 = y + 1 = -\dfrac{z}{2} \)
        \task 法平面方程为 \( x + y - 2z = 0 \)
      \end{tasks}

    \item 函数 \( f(x, y) = \sqrt{x^2 + y^2} \) 在点 \( (0, 0) \) 处（\qquad）.

      \begin{tasks}[label=\textbf{\Alph*.}, label-width=1.5em, column-sep=1em](2)
        \task 可微
        \task 偏导数存在
        \task 连续
        \task 不连续
      \end{tasks}

    \item 已知函数 \( f \) 连续，则二次积分 \( \displaystyle \int_0^{\frac{\pi}{2}} \mathrm{d}\theta \displaystyle \int_0^{2\sin\theta} f(r\cos\theta, r\sin\theta) r\,\mathrm{d}r = \)（\qquad）.

      \begin{tasks}[label=\textbf{\Alph*.}, label-width=1.5em, column-sep=1em](1)
        \task \( \displaystyle \int_0^2 \mathrm{d}x \displaystyle \int_0^{1 + \sqrt{1 - x^2}} f(x, y) \mathrm{d}y \)
        \task \( \displaystyle \int_0^2 \mathrm{d}x \displaystyle \int_0^{\sqrt{2x - x^2}} f(x, y) \mathrm{d}y \)
        \task \( \displaystyle \int_0^2 \mathrm{d}y \displaystyle \int_0^{\sqrt{2y - y^2}} f(x, y) \mathrm{d}x \)
        \task \( \displaystyle \int_0^2 \mathrm{d}x \displaystyle \int_0^2 f(x, y) \mathrm{d}y \)
      \end{tasks}

    \item 设 \( \Gamma: \begin{cases} x^2 + y^2 + z^2 = 1 \\ x + y + z = 0 \end{cases} \)，\( I = \displaystyle \oint_{\Gamma} x\,\mathrm{d}s \)，\( J = \displaystyle \oint_{\Gamma} y\,\mathrm{d}s \)，\( K = \displaystyle \oint_{\Gamma} z^2 \mathrm{d}s \).以下说法中正确的是（\qquad）.

      \begin{tasks}[label=\textbf{\Alph*.}, label-width=1.5em, column-sep=1em](2)
        \task \( K = 0 \)
        \task \( I, J, K \) 中有两个等于 0
        \task \( I, J, K \) 都等于 0
        \task \( I, J, K \) 全都不等于 0
      \end{tasks}

    \item 设 \( f(x) \) 是以 \( 2\pi \) 为周期的周期函数且 \( f(x) = \begin{cases} 0, & -\pi \leq x < 0, \\ \pi - x, & 0 \leq x < \pi \end{cases} \)，\( f(x) \) 的傅里叶级数的和函数是 \( S(x) \)。以下说法中正确的是（\qquad）.

      \begin{tasks}[label=\textbf{\Alph*.}, label-width=1.5em, column-sep=1em](2)
        \task \( S(x) \) 处处连续
        \task \( S(x) \equiv f(x) \)
        \task \( S(-1) = 0 \)
        \task \( S(0) = \pi \)
      \end{tasks}
  \end{enumerate}

  \section{填空题(每小题 4 分，共 16 分)}
  \begin{enumerate}
    \setcounter{enumi}{6}
    \item 直线 \( \dfrac{x-2}{2} = \dfrac{y-1}{1} = \dfrac{z-3}{1} \) 与平面 \( x - y + 2z + 4 = 0 \) 的夹角是 \(\underline{\hspace{5em}}\).
      \vspace{1em}

    \item 设 \( P_0(1,1,-1) \)，\( P_1(2,-1,0) \)，则 \( u = x + y^2 + z^3 \) 在 \( P_0 \) 处沿 \( \overrightarrow{P_0P_1} \) 方向的方向导数为 \(\underline{\hspace{5em}}\).
      \vspace{1em}

    \item 若 \( z = z(x,y) \) 由方程 \( e^{x+2y+3z} + xyz = 1 \) 确定，则 \( z_x(0,0) = \underline{\hspace{5em}}\).
      \vspace{1em}

    \item 级数 \( \sum \limits_{n=0}^{\infty} \dfrac{(-1)^n x^{2n}}{(2n)!} \) (\( -\infty < x < \infty \)) 的和函数为 \(\underline{\quad\quad}\).
      \vspace{1em}
  \end{enumerate}

  \section{基本计算题(每小题 7 分，共 42 分)}

  \begin{enumerate}
    \setcounter{enumi}{10}
    \item 求微分方程 \( x \dfrac{\mathrm{d}y}{\mathrm{d}x} = y \ln \dfrac{y}{x} \) 的通解.
      \vspace{10em}

    \item 已知函数 \( z = f(xy, yg(x)) \)，其中 \( f \) 有二阶连续偏导数，\( g \) 可导，求 \( z_x \)，\( z_{xy} \).
      \vspace{10em}

    \item 计算二次积分 \( I = \displaystyle \int_1^3 \mathrm{d}x \displaystyle \int_{x-1}^2 \sin y^2 \mathrm{d}y \).
      \vspace{12em}

    \item 求三重积分 \( I = \displaystyle \iiint \limits_V (x^3 + y^2 + z) dv \)，其中 \( V \) 为 \( x^2 + y^2 + z^2 \leq a^2 \)，\( a > 0 \).
      \vspace{12em}

    \item 求 \( I = \displaystyle \iint \limits_S (z^2 + x) \mathrm{d}y\,\mathrm{d}z - z\,\mathrm{d}x\,\mathrm{d}y \)，其中 \( S \) 是 \( z = \dfrac{1}{2}(x^2 + y^2) \) (\( 0 \leq z \leq 2 \)) 下侧.
      \vspace{12em}

    \item 将 \( f(x) = \arctan x \) 展开为 Maclaurin 级数，并求 \( f^{(20)}(0) \)，\( f^{(21)}(0) \).
      \vspace{12em}
  \end{enumerate}

  \section{应用题(每小题 7 分，共 14 分)}

  \begin{enumerate}
    \setcounter{enumi}{16}
    \item 求 \( f(x, y) = x^4 + y^4 - x^2 - 2xy - y^2 \) 的极值.
      \vspace{10em}

    \item 已知曲线积分 \( \displaystyle \int_L y f(x) \mathrm{d}x + [f(x) - x^2] \mathrm{d}y \) 与路径无关，其中 \( f(x) \) 有一阶连续导数，且 \( f(0) = 1 \)，求 \( f(x) \) 和 \( I = \displaystyle \int_{(0,0)}^{(1,1)} y f(x) dx + [f(x) - x^2] dy \) 的值.
      \vspace{10em}
  \end{enumerate}

  \section{综合题(每小题 5 分，共 10 分)}

  \begin{enumerate}
    \setcounter{enumi}{18}
    \item 设 \( D: x^2 + y^2 \leq 1 \)，证明不等式：\( \displaystyle \iint \limits_D \sin \sqrt{(x^2 + y^2)^3} \mathrm{d}x\,\mathrm{d}y \leq \dfrac{2}{5}\pi \).
      \vspace{10em}

    \item 设 \( f(0) = 1 \)，\( f'(0) = 0 \)，\( f''(x) \) 在 \( (-1,1) \) 内有界，证明：\( \alpha > \dfrac{1}{2} \) 时，\( \sum \limits_{n=1}^{\infty} \left( f\left( \dfrac{1}{n^\alpha} \right) - 1 \right) \) 绝对收敛.
      \vspace{10em}
  \end{enumerate}
\chapter{2020-2021学年微积分（一）（下）期末考试参考答案}
  \section{单项选择题(每小题 3 分，共 18 分)}
  \begin{enumerate}
    \item \textbf{Solution}. D.

      右端 \(2x-2e^x\) 可分成一次多项式项与指数项两部分. 齐次方程特征方程为
      \[
        r^2-3r+2=(r-1)(r-2)=0,
      \]
      故 \(e^x\) 已在齐次解中出现，因此对应特解应乘以 \(x\). 所以可设
      \[
        y^*=Ax+B+Cxe^x.
      \]

    \item \textbf{Solution}. B.

      两曲面的法向量分别为
      \[
        \nabla(x^2+y^2+z^2)=(2x,2y,2z),\quad \nabla(x+y+z)=(1,1,1).
      \]
      在 \(M(1,-1,0)\) 处取值得
      \[
        (2,-2,0),\quad (1,1,1).
      \]
      叉积可取为
      \[
        (2,-2,0)\times(1,1,1)=(-2,-2,4),
      \]
      所以 A、C、D 都正确，而 B 不是切向方向.

    \item \textbf{Solution}. C.

      函数
      \[
        f(x,y)=\sqrt{x^2+y^2}
      \]
      在 \((0,0)\) 连续，但沿 \(x\) 轴有
      \[
        \frac{f(h,0)-f(0,0)}{h}=\frac{|h|}{h},
      \]
      极限不存在，所以偏导数不存在，更不可微.

    \item \textbf{Solution}. C.

      极坐标区域为
      \[
        0\le \theta\le \frac{\pi}{2},\quad 0\le r\le 2\sin\theta,
      \]
      即第一象限内圆
      \[
        x^2+(y-1)^2\le 1
      \]
      的部分. 改写成直角坐标可得
      \[
        0\le y\le 2,\quad 0\le x\le \sqrt{2y-y^2},
      \]
      故选 C.

    \item \textbf{Solution}. B.

      由对称性可知
      \[
        \oint_\Gamma x\,\mathrm{d}s=\oint_\Gamma y\,\mathrm{d}s=0.
      \]
      但 \(z^2\ge 0\) 且不恒为零，因此
      \[
        \oint_\Gamma z^2\,\mathrm{d}s>0.
      \]
      所以恰有两个积分为零.

    \item \textbf{Solution}. C.

      傅里叶级数在连续点收敛到原函数值，在跳跃点收敛到左右极限平均值. 因为 \(-1\in(-\pi,0)\)，故
      \[
        S(-1)=f(-1)=0.
      \]
      而
      \[
        S(0)=\frac{0+\pi}{2}=\frac{\pi}{2}\neq \pi.
      \]
  \end{enumerate}

  \section{填空题(每小题 4 分，共 16 分)}
  \begin{enumerate}
    \setcounter{enumi}{6}
    \item \textbf{Solution}. $\dfrac{\pi}{6}$.

      直线方向向量为 \((2,1,1)\)，平面法向量为 \((1,-1,2)\). 设夹角为 \(\theta\)，则
      \[
        \sin\theta=\frac{|(2,1,1)\cdot(1,-1,2)|}{\sqrt6\cdot\sqrt6}=\frac12.
      \]
      所以
      \[
        \theta=\frac{\pi}{6}.
      \]

    \item \textbf{Solution}. $0$.

      \[
        \nabla u=(1,2y,3z^2),
      \]
      在 \(P_0(1,1,-1)\) 处为 \((1,2,3)\). 又
      \[
        \overrightarrow{P_0P_1}=(1,-2,1),
      \]
      二者点积为 \(0\)，故方向导数为 \(0\).

    \item \textbf{Solution}. $-\dfrac13$.

      对
      \[
        e^{x+2y+3z}+xyz=1
      \]
      关于 \(x\) 求偏导，得
      \[
        e^{x+2y+3z}(1+3z_x)+yz+xz\,z_x=0.
      \]
      在 \((0,0)\) 处有 \(z=0\)，从而
      \[
        1+3z_x(0,0)=0,\qquad z_x(0,0)=-\frac13.
      \]

    \item \textbf{Solution}. $\cos x$.

      这是余弦函数的 Maclaurin 级数：
      \[
        \cos x=\sum_{n=0}^{\infty}\frac{(-1)^n x^{2n}}{(2n)!}.
      \]
  \end{enumerate}

  \section{基本计算题(每小题 7 分，共 42 分)}
  \begin{enumerate}
    \setcounter{enumi}{10}
    \item \textbf{Solution}. 令
      \[
        u=\frac{y}{x},\qquad y=ux.
      \]
      则
      \[
        \frac{\mathrm{d}y}{\mathrm{d}x}=u+x\frac{\mathrm{d}u}{\mathrm{d}x}.
      \]
      原方程化为
      \[
        x\left(u+x\frac{\mathrm{d}u}{\mathrm{d}x}\right)=ux\ln u,
      \]
      即
      \[
        x\frac{\mathrm{d}u}{\mathrm{d}x}=u(\ln u-1).
      \]
      分离变量并积分得
      \[
        \frac{\mathrm{d}u}{u(\ln u-1)}=\frac{\mathrm{d}x}{x}
        \Longrightarrow
        \ln\frac{y}{x}=Cx+1.
      \]

    \item \textbf{Solution}. 记
      \[
        u=xy,\qquad v=yg(x),
      \]
      则 \(z=f(u,v)\). 由链式法则
      \[
        z_x=f_1'u_x+f_2'v_x=y(f_1'+g'(x)f_2').
      \]
      再对 \(y\) 求偏导，得
      \[
        z_{xy}=f_1'+g'(x)f_2'+xyf_{11}''+[g(x)+xyg'(x)]f_{12}''+yg'(x)g(x)f_{22}''.
      \]

    \item \textbf{Solution}. 原积分区域为
      \[
        1\le x\le 3,\quad x-1\le y\le 2.
      \]
      交换积分次序后可写成
      \[
        0\le y\le 2,\quad 1\le x\le y+1.
      \]
      因而
      \[
        I=\int_0^2\,\mathrm{d}y\int_1^{y+1}\sin y^2\,\mathrm{d}x
        =\int_0^2 y\sin y^2\,\mathrm{d}y.
      \]
      令 \(t=y^2\)，则
      \[
        I=\frac12\int_0^4\sin t\,\mathrm{d}t
        =\frac12(1-\cos4).
      \]

    \item \textbf{Solution}. 由球域关于原点的对称性，
      \[
        \iiint_V x^3\,\mathrm{d}v=0,\qquad \iiint_V z\,\mathrm{d}v=0,
      \]
      所以
      \[
        I=\iiint_V y^2\,\mathrm{d}v.
      \]
      由轮换对称性，
      \[
        3I=\iiint_V(x^2+y^2+z^2)\,\mathrm{d}v.
      \]
      用球坐标计算得
      \[
        3I=\int_0^{2\pi}\int_0^\pi\int_0^a \rho^4\sin\varphi\,\mathrm{d}\rho\,\mathrm{d}\varphi\,\mathrm{d}\theta
        =\frac{4\pi a^5}{5},
      \]
      所以
      \[
        I=\frac{4\pi a^5}{15}.
      \]

    \item \textbf{Solution}. 记
      \[
        P=z^2+x,\qquad Q=0,\qquad R=-z.
      \]
      将侧面与上盖平面 \(z=2\) 补成闭曲面，所围区域记为 \(\Omega\). 由高斯公式
      \[
        \iint_{\partial\Omega}P\,\mathrm{d}y\,\mathrm{d}z+Q\,\mathrm{d}z\,\mathrm{d}x+R\,\mathrm{d}x\,\mathrm{d}y
        =\iiint_\Omega(P_x+Q_y+R_z)\,\mathrm{d}v.
      \]
      这里
      \[
        P_x+Q_y+R_z=1-1=0,
      \]
      所以侧面与上盖面的通量之和为零. 再计算上盖面的贡献，可得原积分
      \[
        I=8\pi.
      \]

    \item \textbf{Solution}. 由
      \[
        \frac{1}{1+x^2}=\sum_{n=0}^{\infty}(-1)^n x^{2n}\qquad (|x|<1)
      \]
      两端积分，得
      \[
        \arctan x=\sum_{n=0}^{\infty}\frac{(-1)^n}{2n+1}x^{2n+1},\qquad |x|\le 1.
      \]
      比较 Maclaurin 展开系数可知
      \[
        f^{(20)}(0)=0,\qquad f^{(21)}(0)=20!.
      \]
  \end{enumerate}

  \section{应用题(每小题 7 分，共 14 分)}
  \begin{enumerate}
    \setcounter{enumi}{16}
    \item \textbf{Solution}. 由
      \[
        f_x=4x^3-2x-2y,\qquad f_y=4y^3-2x-2y
      \]
      解驻点方程可得
      \[
        (x,y)=(0,0),(1,1),(-1,-1).
      \]
      再由
      \[
        f_{xx}=12x^2-2,\quad f_{yy}=12y^2-2,\quad f_{xy}=-2
      \]
      判别知 \((1,1)\) 与 \((-1,-1)\) 是极小值点，且
      \[
        f(1,1)=f(-1,-1)=-2.
      \]
      点 \((0,0)\) 不是极值点.

    \item \textbf{Solution}. 记
      \[
        P(x,y)=yf(x),\qquad Q(x,y)=f(x)-x^2.
      \]
      由路径无关得
      \[
        P_y=Q_x,
      \]
      即
      \[
        f'(x)-f(x)=2x.
      \]
      解得
      \[
        f(x)=Ce^x-2x-2.
      \]
      再由 \(f(0)=1\) 得 \(C=3\)，故
      \[
        f(x)=3e^x-2x-2.
      \]
      取折线路径 \((0,0)\to(1,0)\to(1,1)\)，第一段积分为 \(0\)，第二段积分为
      \[
        \int_0^1[f(1)-1]\mathrm{d}y=f(1)-1=3e-5.
      \]
      所以
      \[
        I=3e-5.
      \]
  \end{enumerate}

  \section{综合题(每小题 5 分，共 10 分)}
  \begin{enumerate}
    \setcounter{enumi}{18}
    \item \textbf{Proof}. 令
      \[
        x=r\cos\theta,\qquad y=r\sin\theta,
      \]
      则
      \[
        \iint_D \sin \sqrt{(x^2+y^2)^3}\,\mathrm{d}x\,\mathrm{d}y
        =\int_0^{2\pi}\int_0^1 \sin(r^3)\,r\,\mathrm{d}r\,\mathrm{d}\theta
        =2\pi\int_0^1 r\sin(r^3)\,\mathrm{d}r.
      \]
      由 \(\sin t\le t\)（\(t\ge 0\)）得
      \[
        2\pi\int_0^1 r\sin(r^3)\,\mathrm{d}r
        \le 2\pi\int_0^1 r\cdot r^3\,\mathrm{d}r
        =2\pi\int_0^1 r^4\,\mathrm{d}r
        =\frac{2\pi}{5}.
      \]
      于是原不等式成立.

    \item \textbf{Proof}. 因 \(f''(x)\) 在 \((-1,1)\) 内有界，设
      \[
        |f''(x)|\le M\qquad (|x|<1).
      \]
      对 \(x_n=\dfrac{1}{n^\alpha}\)，由 Taylor 公式
      \[
        f(x_n)=f(0)+f'(0)x_n+\frac{f''(\theta_n x_n)}{2}x_n^2,
      \]
      其中 \(0<\theta_n<1\). 利用 \(f(0)=1,\ f'(0)=0\) 得
      \[
        f\left(\frac{1}{n^\alpha}\right)-1
        =\frac{f''(\theta_n/n^\alpha)}{2n^{2\alpha}}.
      \]
      所以
      \[
        \left|f\left(\frac{1}{n^\alpha}\right)-1\right|
        \le \frac{M}{2n^{2\alpha}}.
      \]
      当 \(\alpha>\dfrac12\) 时，级数
      \[
        \sum_{n=1}^{\infty}\frac{1}{n^{2\alpha}}
      \]
      收敛，于是由比较判别法知
      \[
        \sum_{n=1}^{\infty}\left|f\left(\frac{1}{n^\alpha}\right)-1\right|
      \]
      收敛，即原级数绝对收敛.
  \end{enumerate}
\end{document}






